\begin{problem}[第34届物理决赛][毛细管水滴]{85}
    有一根长为 $6.00\,\mathrm{cm}$、内外半径分别为 $0.500\,\mathrm{mm}$ 和 $5.00\,\mathrm{mm}$ 的玻璃毛细管。
    \begin{subproblem}
        \begin{subsubproblem}
            毛细管竖直悬空固定放置，注入水后，在管的下端中央形成一悬挂的水滴，管中水柱表面中心相对于水滴底部的高度为 $3.50\,\mathrm{cm}$，求水滴底部表面的曲率半径 $a$；
        \end{subsubproblem}
        \begin{subsubproblem}
            若将该毛细管长度的三分之一竖直浸入水中，问需要多大向上的力才能使该毛细管保持不动？
            已知玻璃的密度是水的 $2$ 倍，水的密度为 $1.00\times10^{3}\,\mathrm{kg\cdot m^{-3}}$，水的表面张力系数为 $7.27\times10^{-2}\,\mathrm{N\cdot m^{-1}}$，水与玻璃的接触角 $\theta$ 可视为零，重力加速度取 $9.80\,\mathrm{m\cdot s^{-2}}$。
        \end{subsubproblem}
    \end{subproblem}
\end{problem}